% Footnote with ATLAS coordinate system
\newcommand{\AtlasCoordFootnote}{%
ATLAS uses a right-handed coordinate system with its origin at the nominal interaction point (IP)
in the centre of the detector, the \(z\)-axis along the beam pipe, the \(x\)-axis pointing to the centre of the LHC ring,
and the \(y\)-axis pointing upwards.
Cylindrical coordinates \((r,\phi)\) are used in the transverse plane,
\(\phi\) being the azimuthal angle around the \(z\)-axis.
The pseudorapidity is defined in terms of the polar angle \(\theta\) as \(\eta = -\ln \tan(\theta/2)\).
Angular distances are measured by \(\Delta R \equiv \sqrt{(\Delta\eta)^{2} + (\Delta\phi)^{2}}\).}

%-------------------------------------------------------------------------------
The ATLAS detector~\cite{PERF-2007-01} at the LHC consists of an inner tracking detector system (ID) surrounded by a thin superconducting solenoid, an electromagnetic and a hadron calorimeters (ECAL and HCAL),
and a muon spectrometer (MS) incorporating three large superconducting air-core toroidal magnets, with 
an almost complete solid angle coverage around the collision point.
%\footnote{\AtlasCoordFootnote}.
\AtlasCoordFootnote

The ID system is immersed in a \SI{2}{\tesla} axial magnetic field 
and provides charged-particle tracking in the range \(|\eta| < 2.5\).
The high-granularity silicon pixel detector covers the vertex region and typically provides four measurements per track, 
the first hit normally being in the insertable B-layer installed before Run~2~\cite{ATLAS-TDR-19,PIX-2018-001}.
It is followed by the silicon microstrip tracker, which usually provides eight measurements per track.
The outermost part of the ID is the transition radiation tracker (TRT),
which enables radially extended track reconstruction up to \(|\eta| = 2.0\). 
The TRT also provides electron identification information 
based on the fraction of hits (typically 30 in total) above a higher energy-deposit threshold corresponding to transition radiation.

The calorimeter system covers the pseudorapidity range \(|\eta| < 4.9\). 
The ECAL~\cite{ATLAS-TDR-02,LARG-2009-01} is a lead/liquid-argon (LAr) sampling calorimeter covering the region \(|\eta|< 3.2\), divided into a barrel and two endcaps, and longitudinally segmented into three layers: the strips, the middle and the back layer. The 3-dimensional high-granularity in the fiducial acceptances $|\eta|<1.37$ and $1.52<|\eta|<2.37$ allows to observe the shape of the electromagnetic showers. Before the ECAL, a thin LAr presampler covering \(|\eta| < 1.8\) detects energy losses of showers starting in the material upstream.
% electromagnetic (EM) calorimetry is provided by barrel and 
% endcap high-granularity lead/liquid-argon (LAr) calorimeters,
% with an additional thin LAr presampler covering \(|\eta| < 1.8\)
% to correct for energy loss in material upstream of the calorimeters.
The HCAL is made by the steel/scintillator-tile calorimeter,
segmented into three barrel structures within \(|\eta| < 1.7\), and two copper/LAr hadron endcap calorimeters.
The solid angle coverage is completed with forward copper/LAr and tungsten/LAr calorimeter modules
optimised for electromagnetic and hadronic energy measurements respectively.

The MS comprises separate trigger and
high-precision tracking chambers measuring the deflection of muons in a magnetic field generated by the superconducting air-core toroidal magnets.
The field integral of the toroids ranges between \num{2.0} and \SI{6.0}{\tesla\metre}
across most of the detector. 
Three layers of precision chambers, each consisting of layers of monitored drift tubes, covers the region \(|\eta| < 2.7\),
complemented by cathode-strip chambers in the forward region, where the background is highest.
The muon trigger system covers the range \(|\eta| < 2.4\) with resistive-plate chambers in the barrel, and thin-gap chambers in the endcap regions.

Interesting events are selected by the first-level (L1) trigger system implemented in custom hardware,
followed by selections made by software-implemented algorithms in the high-level trigger (HLT)~\cite{TRIG-2016-01}. 
The first-level trigger selects events at a rate below \SI{100}{\kHz} from the \SI{40}{\MHz} bunch crossings; the high-level trigger further reduces that rate, in order to record events to disk at about \SI{3}{\kHz}.
An extensive software suite~\cite{ATL-SOFT-PUB-2021-001} is used in the reconstruction and analysis of real
and simulated data, in detector operations, and in the trigger and data acquisition systems of the experiment.
